% !PRINTED 311
\weberpage{311}
%% [page furniture, not body text -- commented out 2026-08-10] Page furniture (running
%% head, folio, signature) is excluded from this edition by standing policy.
%% \noindent \S.\,81. \hfill Triple systems.

%<<CONT
are satisfied for three and only three systems of the unknowns $\xi_1, \xi_2, \xi_3$, and that these three form a triple, so that the triple is determined by the three congruences
\weq{5}{\Sigma \alpha\xi \equiv 0, \quad \Sigma \alpha\eta \equiv 0, \quad \Sigma \alpha\zeta \equiv 0 \quad (\mathrm{mod}\ 2)}
One recognises this without lengthy general considerations if one considers individually the cases in which only one $\alpha$, or two, or all three $\alpha \equiv 1$ (mod 2).

Since there are seven distinct systems of numbers $\alpha_1, \alpha_2, \alpha_3$, a triple is completely determined by the relations (5) from each such system of the $\alpha$, and we can denote the function $v$ unambiguously by $v_{\alpha_1,\alpha_2,\alpha_3}$ or more briefly by $v_\alpha$:
\weq{6}{v_\alpha = \frac{1}{\Psi_\xi \Psi_\eta \Psi_\zeta}.}

\gesperrt{3. The functions $v_\alpha$, regarded as functions of the $X$, admit the permutations of the group $Q$.}

If we apply to the $z$ a substitution $\sigma$ from $S$, then by 2.\ (\S.\,80) the $\xi, \eta, \zeta$ simultaneously undergo the substitution $\varrho^{-1}$; from this, however, it follows by (5) that the $\alpha$ undergo the substitution $\sigma$, and we therefore have:

\gesperrt{4. The application of the substitution $\sigma$ from $S$ to the indices $z$ of $X$ has as its consequence the application of the same substitution to the indices $\alpha$ of $v_\alpha$.}

If, in the sums (5), in which $(\xi), (\eta), (\zeta)$ is the triple belonging to $(\alpha)$, we put for $(\alpha)$ another system $(\beta)$, then on account of (3)
\[
\Sigma \beta\xi + \Sigma \beta\eta + \Sigma \beta\zeta \equiv 0.
\]

Now if $\beta$ is different from $\alpha$, these three sums cannot all three be congruent to 0, because otherwise two distinct systems $(\alpha), (\beta)$ would lead to the same triple; \gesperrt{consequently, of these sums, two must be odd and one even.} We therefore have the theorem:

\gesperrt{5. If $(\xi)$ runs through the triple belonging to $(\alpha)$, and $(\beta)$ is a system different from $(\alpha)$, then among the sums $\Sigma \beta\xi$ two are odd and one is even.}
